% GENERATED FILE. Edit canonical lessons, then run npm run generate:books.
% canonical-source-hash: fnv1a64:0f4d6919545efffc
% canonical-lessons: JA-C110-hae, JA-C110-katatsumuri, JA-C110-shippo, JA-C110-kakato, JA-C110-sune

\chapter{Fly, Snail, Tail, Heel, Shin}
\label{ch:ja-fly-snail-tail-heel-shin}
\hlchaptermodality{110}

\begin{chapteropening}
\textbf{By the end of this chapter:} \emph{I can say a fly, a snail, a tail, a heel and a shin.}

Complete the last lesson of chapter 110: I can say a fly, a snail, a tail, a heel and a shin.
\end{chapteropening}

\section[\texorpdfstring{\ja{はえ} (hae)}{はえ (hae)}]{\ja{はえ} (hae) — a fly}

\label{lesson:JA-C110-hae}



\begin{quote}
Before the new one: say the Japanese for an owl, then the Japanese for a silkworm.
\end{quote}

\subsection*{You'll want to know: \ja{はえ}}
\textbf{\ja{はえ}} — \emph{hae} — \textquotedblleft{}a fly\textquotedblright{}.

\begin{grammarlens}[title={What you've built}]
One more word for this chapter.
\end{grammarlens}

\subsection*{Your turn}
\begin{itemize}
\raggedright
  \item \emph{Say it:} \emph{hae}
  \item \emph{Say it:} \emph{hae}, once more
  \item \emph{Say it:} say \emph{hae}
  \item \emph{Recall:} say the Japanese for an owl, then the Japanese for a silkworm, then say \emph{hae} again
\end{itemize}

\begin{tcolorbox}[breakable,colback=teal!4,colframe=teal!35!black,title={Before you move on}]
What does \ja{はえ} mean? (\textquotedblleft{}A fly\textquotedblright{}.) Say it once more.
\end{tcolorbox}
\section[\texorpdfstring{\ja{かたつむり} (katatsumuri)}{かたつむり (katatsumuri)}]{\ja{かたつむり} (katatsumuri) — a snail}

\label{lesson:JA-C110-katatsumuri}



\begin{quote}
Before the new one: say the Japanese for a silkworm, then the Japanese for a fly.
\end{quote}

\subsection*{You'll want to know: \ja{かたつむり}}
\textbf{\ja{かたつむり}} — \emph{katatsumuri} — \textquotedblleft{}a snail\textquotedblright{}.

\begin{grammarlens}[title={What you've built}]
One more word for this chapter.
\end{grammarlens}

\subsection*{Your turn}
\begin{itemize}
\raggedright
  \item \emph{Say it:} \emph{katatsumuri}
  \item \emph{Say it:} \emph{katatsumuri}, once more
  \item \emph{Say it:} say \emph{katatsumuri}
  \item \emph{Recall:} say the Japanese for a silkworm, then the Japanese for a fly, then say \emph{katatsumuri} again
\end{itemize}

\begin{tcolorbox}[breakable,colback=teal!4,colframe=teal!35!black,title={Before you move on}]
What does \ja{かたつむり} mean? (\textquotedblleft{}A snail\textquotedblright{}.) Say it once more.
\end{tcolorbox}
\section[\texorpdfstring{\ja{しっぽ} (shippo)}{しっぽ (shippo)}]{\ja{しっぽ} (shippo) — a tail}

\label{lesson:JA-C110-shippo}



\begin{quote}
Before the new one: say the Japanese for a fly, then the Japanese for a snail.
\end{quote}

\subsection*{You'll want to know: \ja{しっぽ}}
\textbf{\ja{しっぽ}} — \emph{shippo} — \textquotedblleft{}a tail\textquotedblright{}.

\begin{grammarlens}[title={What you've built}]
One more word for this chapter.
\end{grammarlens}

\subsection*{Your turn}
\begin{itemize}
\raggedright
  \item \emph{Say it:} \emph{shippo}
  \item \emph{Say it:} \emph{shippo}, once more
  \item \emph{Say it:} say \emph{shippo}
  \item \emph{Recall:} say the Japanese for a fly, then the Japanese for a snail, then say \emph{shippo} again
\end{itemize}

\begin{tcolorbox}[breakable,colback=teal!4,colframe=teal!35!black,title={Before you move on}]
What does \ja{しっぽ} mean? (\textquotedblleft{}A tail\textquotedblright{}.) Say it once more.
\end{tcolorbox}
\section[\texorpdfstring{\ja{かかと} (kakato)}{かかと (kakato)}]{\ja{かかと} (kakato) — a heel}

\label{lesson:JA-C110-kakato}



\begin{quote}
Before the new one: say the Japanese for a snail, then the Japanese for a tail.
\end{quote}

\subsection*{You'll want to know: \ja{かかと}}
\textbf{\ja{かかと}} — \emph{kakato} — \textquotedblleft{}a heel\textquotedblright{}.

\begin{grammarlens}[title={What you've built}]
One more word for this chapter.
\end{grammarlens}

\subsection*{Your turn}
\begin{itemize}
\raggedright
  \item \emph{Say it:} \emph{kakato}
  \item \emph{Say it:} \emph{kakato}, once more
  \item \emph{Say it:} say \emph{kakato}
  \item \emph{Recall:} say the Japanese for a snail, then the Japanese for a tail, then say \emph{kakato} again
\end{itemize}

\begin{tcolorbox}[breakable,colback=teal!4,colframe=teal!35!black,title={Before you move on}]
What does \ja{かかと} mean? (\textquotedblleft{}A heel\textquotedblright{}.) Say it once more.
\end{tcolorbox}
\section[\texorpdfstring{\ja{すね} (sune)}{すね (sune)}]{\ja{すね} (sune) — a shin}

\label{lesson:JA-C110-sune}



\begin{quote}
Before the new one: say the Japanese for a tail, then the Japanese for a heel.
\end{quote}

\subsection*{You'll want to know: \ja{すね}}
\textbf{\ja{すね}} — \emph{sune} — \textquotedblleft{}a shin\textquotedblright{}.

\begin{grammarlens}[title={What you've built}]
One more word for this chapter.
\end{grammarlens}

\subsection*{Your turn}
\begin{itemize}
\raggedright
  \item \emph{Say it:} \emph{sune}
  \item \emph{Say it:} \emph{sune}, once more
  \item \emph{Say it:} say \emph{sune}
  \item \emph{Recall:} say the Japanese for a tail, then the Japanese for a heel, then say \emph{sune} again
\end{itemize}

\begin{tcolorbox}[breakable,colback=teal!4,colframe=teal!35!black,title={Before you move on}]
What does \ja{すね} mean? (\textquotedblleft{}A shin\textquotedblright{}.) Say it once more.
\end{tcolorbox}
